%! AP Statistics — Unit 3, Lesson 3.5 — Activity KEY
\documentclass[10pt]{article}
\usepackage{apstats-article}
\usepackage{apstats-key}

\begin{document}

\pageheader{Unit 3, Lesson 3.5}{Group Activity KEY: Designing Experiments}
\namepartnerperiod

\vspace{0.1in}

\begin{scenariobox}[Context]{sky}
\small
Your group will analyze and design experiments. For each scenario, identify the
components, evaluate whether the four principles are met, and improve the design
where needed. Then design your own experiment from scratch.
\end{scenariobox}

\vspace{0.1in}

\textbf{Scenario A:} A researcher wants to know if a new supplement improves memory.
She recruits 30 volunteers. Those who ask to take the supplement get it; the rest get
nothing. After 4 weeks, all participants take a memory test.

\begin{practicebox}
\small
\begin{enumerate}[leftmargin=1.5em, itemsep=7pt]
  \item Experimental units: \makebox[8cm][l]{\ans{The 30 volunteers}}
  \item Factor: \makebox[8cm][l]{\ans{Supplement (supplement vs.\ no supplement)}}
  \item Treatments: \makebox[8cm][l]{\ans{Take the supplement; take nothing}}
  \item Response variable: \makebox[8cm][l]{\ans{Memory test score}}
  \item Which of the four principles is \emph{violated}? \makebox[5cm][l]{\ans{Random assignment}} Why?\\
        \ansline{Participants chose their own treatment, so the groups may}\\[1pt]\ansline{differ in ways other than the supplement.}
  \item What design flaw allows a confounding variable? Describe it.\\
        \ansline{Self-selection: motivated volunteers may already sleep or study}\\[1pt]\ansline{more, which improves memory regardless of the supplement.}
  \item Redesign: How would you fix this experiment?\\
        \ansline{Randomly assign the 30 to supplement or a placebo pill, and}\\[1pt]\ansline{keep subjects (and the evaluator) blind to the group.}
\end{enumerate}
\end{practicebox}

\vspace{0.1in}

\textbf{Scenario B:} Researchers test two brands of headache medication. They have
40 patients with chronic headaches. Half are randomly assigned to Brand A;
half to Brand B. Neither patients nor doctors know which brand a patient is taking.
After 2 weeks, doctors rate each patient's improvement.

\begin{practicebox}
\small
\begin{enumerate}[leftmargin=1.5em, itemsep=7pt]
  \item Design type (CRD, block, matched pairs): \makebox[6cm][l]{\ans{CRD}}
  \item What type of blinding is used? \makebox[6cm][l]{\ans{Double-blind}}
  \item Is there a control group? \makebox[3cm][l]{\ans{No.}} Should there be? Explain.\\
        \ansline{Not required to compare two active brands, but a placebo}\\[1pt]\ansline{group would show whether either brand works at all.}
  \item Patients vary in age (some are 20--30 years old; some are 60--70). Would
        blocking by age group improve this design? Explain.\\
        \ansline{Yes --- if age affects response, blocking by age removes that}\\[1pt]\ansline{variability and gives a more precise comparison of brands.}
\end{enumerate}
\end{practicebox}

\vspace{0.1in}

\textbf{Part 3: Design your own experiment.}

\begin{practicebox}
\small
Research question: \textit{Does listening to classical music while studying improve
quiz scores for Statistics students?}

\vspace{6pt}
Using 30 Statistics students, design a well-controlled experiment. Address:
\begin{enumerate}[leftmargin=1.5em, itemsep=6pt]
  \item Treatments (be specific): \makebox[9cm][l]{\ans{Study with classical music vs.\ study in silence.}}
  \item How you will randomly assign students: \makebox[9cm][l]{\ans{Number 1--30; randInt picks 15 for music.}}
  \item Control of confounding variables (list 2 and how you address each):\\
        1. \makebox[9cm][l]{\ans{Study material: same topic and notes for all.}}\\[4pt]
        2. \makebox[9cm][l]{\ans{Setting: same room, same time, same study length.}}
  \item Response variable and how it will be measured: \makebox[9cm][l]{\ans{Score on the same quiz right after studying.}}
  \item Will you use blinding? What type is feasible? \makebox[9cm][l]{\ans{Single-blind: the grader doesn't know groups.}}
  \item \makebox[0pt][l]{\raisebox{-2.2\baselineskip}[0pt][0pt]{\parbox[t]{0.9\linewidth}{\ans{30 students $\to$ random assignment $\to$ Group 1 (15): classical music\\
        \phantom{30 students $\to$ random assignment $\to$ }Group 2 (15): silence\\
        $\to$ Compare mean quiz scores.}}}}Draw a diagram of your design in the space below.
  \vspace{1.6in}
\end{enumerate}
\end{practicebox}

\end{document}
